\section{Interpretability: Opening the Box}
\label{sec:interpretability}

Cadence ships a detector you can read like a recipe: gate the leg as moving,
compute the Freeze Index, compare it against a single printed number. Nothing is
hidden. But the project also explored two richer learners --- the FoGNet CNN of the
previous section and, sitting between the two extremes, a small \emph{glass-box}
RandomForest built from nine hand-named features. This section uses that forest as a
worked example of \emph{interpretability}: a set of tools that pry open a trained
model and ask \emph{which inputs it actually leans on, and in which direction}. We
then return to the design question that matters clinically --- why a wearable that
warns a person about their own falls should prefer a rule it can defend over a
network it cannot.

\begin{keyidea}
A \textbf{black box} (the CNN) maps a raw window straight to $P(\text{freeze})$; its
18\,k weights are accurate but mute --- they cannot tell you \emph{why}. A
\textbf{glass box} reasons over named quantities a clinician already understands
(jerk, dominant frequency, band powers), so every verdict comes with an audit trail.
Cadence's shipped Freeze Index is the most transparent box of all: one ratio, one
threshold, fully traceable by hand.
\end{keyidea}

\subsection{The interpretable alternative: a 9-feature RandomForest}

Instead of feeding the network 256 raw samples per axis, the forest is handed nine
summary numbers computed from the same 4\,s window: the band powers Cadence already
measures (\code{loco\_power}, \code{freeze\_power}, \code{tremor\_power},
\code{total\_power}), the $\FI$ itself (\code{freeze\_index}), the magnitude
statistics \code{mag\_rms} and \code{mag\_range}, the dominant frequency
\code{dom\_freq}, and \code{jerk\_rms} (the RMS of the magnitude's time-derivative ---
how abruptly the leg's acceleration changes). A RandomForest is an ensemble of
decision trees, each a chain of \code{if feature $>$ threshold} splits; every split
is a question you can read aloud. Under nested leave-one-subject-out
cross-validation, at its operating point $t^\star\!\approx\!0.10$ the forest reaches
\textbf{sensitivity 0.75} and \textbf{specificity 0.76}, and as a threshold-free
ranker it scores \textbf{ROC-AUC 0.82} (PR-AUC 0.33). That is a balanced,
honest detector --- and, crucially, one whose decisions we can interrogate with the
three tools below.

\begin{takeaway}
The forest trades a sliver of the CNN's raw ceiling for something the network cannot
offer at any price: every prediction decomposes into contributions from features a
neurologist can name.
\end{takeaway}

\subsection{Permutation importance: break a feature, watch the score fall}

The first question of any model is \emph{which inputs does it actually use?} Permutation
importance answers it by sabotage. Take a trained model and a held-out set; shuffle one
feature's column across all windows --- so that feature now carries pure noise while
every other feature, and the model, stay exactly as they were --- and re-score. If
balanced accuracy collapses, the model was relying on that feature; if nothing happens,
the feature was decoration.

\begin{mathsbox}
For feature $j$, importance is the \emph{drop} in balanced accuracy when its column is
randomly permuted:
\begin{equation}
I_j \;=\; \mathrm{BA}_{\text{intact}} \;-\; \mathbb{E}\!\left[\,\mathrm{BA}_{\text{$j$ shuffled}}\,\right].
\label{eq:perm}
\end{equation}
The expectation is taken over several shuffles (the error bars in
Fig.~\ref{fig:perm} are their spread) and averaged across the held-out subjects, so
$I_j$ reads in the same units as accuracy: ``shuffling \code{jerk\_rms} costs the
forest about 0.05 of balanced accuracy.'' Permuting rather than \emph{dropping} the
column is deliberate --- it preserves that feature's marginal distribution and only
destroys its \emph{link to the label}, isolating the information the model drew from it.
\end{mathsbox}

\begin{figure}[H]
\centering
\includegraphics[width=0.86\linewidth]{perm_importance}
\caption{Permutation importance for the 9-feature forest (pooled LOSO, averaged over
subjects; bars are the mean drop in balanced accuracy, whiskers the spread over
shuffles). Energy-and-motion features --- \code{total\_power}, \code{jerk\_rms},
\code{freeze\_power} --- carry the most weight. Strikingly, the hand-built
\code{freeze\_index} lands \emph{last}: once the forest already has the raw band
powers and jerk, the $\FI$ ratio adds little \emph{marginal} information. The footnote
operating point (sensitivity 0.12, specificity 0.98) is the forest at a naive default
threshold --- a vivid preview of \S\ref{sec:contour} on why the threshold, not the
model, sets the trade-off.}
\label{fig:perm}
\end{figure}

Two lessons fall out of Fig.~\ref{fig:perm}. First, the forest's muscle is in raw
motion energy and abruptness, not in the engineered ratio --- a reminder that a clever
hand-feature can be redundant once a learner sees its raw ingredients. Second, that
\code{freeze\_index} ranks low \emph{inside the forest} is not a verdict against the
$\FI$: importance is always \emph{marginal} and \emph{model-specific}. In Cadence's
shipped detector the $\FI$ is not one voice among nine --- it \emph{is} the decision.

\subsection{SHAP: why \emph{this} window was called a freeze}

Permutation importance is a global headline: it ranks features over the whole
cohort. SHAP (SHapley Additive exPlanations) is the close-up. For a single window it
splits the model's output into a sum of per-feature contributions, each signed: a
positive SHAP value \emph{pushed this window toward freeze}, a negative one
\emph{pulled it back toward calm walking}.

\begin{intuition}
Picture the model's prediction as a bill split fairly among nine diners. SHAP borrows
the Shapley value from cooperative game theory: it asks how much each feature
\emph{adds} to the prediction, averaged over every order in which features could be
revealed. A window with high jerk and a 3--8\,Hz dominant frequency gets a large
positive ``jerk'' contribution and a positive ``frequency'' contribution that
\emph{add up} to a high $P(\text{freeze})$ --- and you can read the receipt line by
line.
\end{intuition}

\begin{figure}[H]
\centering
\includegraphics[width=0.78\linewidth]{shap_bar}
\caption{SHAP global summary: mean $|\text{SHAP}|$ per feature, i.e.\ each feature's
average push on the output magnitude across all windows. \code{jerk\_rms} dominates by
a wide margin --- freezing's defining 3--8\,Hz trembling is, physically, rapid
acceleration reversals, which is exactly what jerk measures --- followed by
\code{dom\_freq} and \code{mag\_range}. The ranking broadly echoes
Fig.~\ref{fig:perm} but is computed by a completely different mechanism (output
attribution rather than accuracy sabotage); their agreement is corroboration, not
coincidence.}
\label{fig:shap}
\end{figure}

Averaging the magnitude of every window's SHAP values yields a global bar chart
(Fig.~\ref{fig:shap}) that two independent methods --- permutation and SHAP --- now
agree puts motion-and-jerk features at the top. Where SHAP earns its keep is the
\emph{per-window} explanation it can also produce: for any flagged second, it states
in plain features why the model raised the alarm, which is precisely the artefact a
clinician needs to either trust or contest a single alert.

\subsection{Partial dependence: which way does the feature push?}

Importance tells us \emph{how much} a feature matters; partial dependence tells us
\emph{in which direction}. Sweep one feature across its range, and at each value
average the model's prediction over every window's real settings of the other eight
features. The resulting curve is the feature's marginal effect on $P(\text{freeze})$,
with the rest of the data marginalised out.

\begin{figure}[H]
\centering
\includegraphics[width=0.96\linewidth]{partial_dependence}
\caption{Partial dependence of $P(\text{freeze})$ on the three leading features.
\textbf{Left:} \code{total\_power} is essentially flat --- the forest uses it for
\emph{gating} (is the leg moving at all?) more than for direction. \textbf{Centre:}
\code{jerk\_rms} climbs monotonically; more abrupt acceleration steadily raises the
freeze probability, exactly matching the physiology of 3--8\,Hz trembling and SHAP's
verdict. \textbf{Right:} \code{freeze\_power} drifts gently \emph{down} --- on its own,
once jerk is accounted for, raw freeze-band energy is not independently incriminating.
The rug marks at the foot of each axis show where the data actually live, a reminder
to distrust the curve out in the sparse tails.}
\label{fig:pdp}
\end{figure}

The monotone \code{jerk\_rms} curve is the satisfying one: it is not merely
\emph{important}, it points the physically correct way. This is the litmus test of a
trustworthy model --- not just that it scores well, but that the reasons it scores
well are the reasons a domain expert would give. A model that achieved the same
accuracy by riding \code{jerk\_rms} \emph{downward} would be a red flag, however good
its numbers.

\subsection{The threshold, not the model, sets the trade-off}
\label{sec:contour}

Interpretability also exposes a trap that pure accuracy hides. Figure~\ref{fig:perm}'s
footnote reports the same forest at sensitivity 0.12, specificity 0.98 --- it catches
almost no freezes --- yet the operating point quoted earlier was 0.75 / 0.76. Nothing
about the \emph{model} changed between those two numbers; only the decision threshold
$t$ did. The freeze class is rare, so the forest's posterior $P(\text{freeze})$ rarely
exceeds the naive $t\!=\!0.5$, and a half-line classifier therefore says ``calm''
almost always --- high specificity, near-blind sensitivity.

\begin{figure}[H]
\centering
\includegraphics[width=0.8\linewidth]{decision_contour}
\caption{The forest's posterior $P(\text{freeze})$ shaded over its two most telling
features, with real windows overlaid (red = freeze, grey = no-freeze). The
\textbf{dashed} contour is the naive $t\!=\!0.5$ boundary: it walls off only a tiny
high-jerk corner, abandoning most red points --- this is the sensitivity-0.12 detector.
The \textbf{solid green} contour is the deployed $t\!\approx\!0.10$ boundary, which
sweeps far lower and reclaims the freeze cluster strung along high \code{jerk\_rms}.
Same model, same plot --- moving one number transforms a timid detector into a
balanced one.}
\label{fig:contour}
\end{figure}

\begin{takeaway}
For a rare, high-cost event, the operating threshold \emph{is} a design decision, not a
default. A missed freeze can mean a fall, so both the forest ($t^\star\!\approx\!0.10$)
and the shipped $\FI$ detector ($t\!=\!1.815$, below the Youden-optimal 2.10) are
deliberately biased toward sensitivity --- and being able to \emph{see} that choice in
Fig.~\ref{fig:contour} is itself an argument for glass boxes.
\end{takeaway}

\subsection{Why Cadence ships the Freeze Index, not a black box}

The forest is interpretable; the $\FI$ detector is interpretable \emph{and}
already on the ankle. Three clinical reasons settle the choice.

\begin{description}[leftmargin=1.6em,style=nextline]
\item[Trust.] A person is being told something about their own falls. ``You froze
because your 3--8\,Hz energy was $1.8\times$ your walking energy'' is a sentence a
patient and a clinician can both check; ``the network's logit crossed zero'' is not.
Every Cadence verdict on the dashboard --- the Freeze Margin meter, the raw $\FI$, the
energy gate --- traces back to one printed equation.
\item[Debuggability.] When the $\FI$ detector misfires we can point at the exact cause
--- the gate let through a still-leg artefact, or the window straddled a freeze
boundary and diluted it --- and fix that line. A misfiring CNN offers 18\,k weights and
a shrug. (Indeed, the design lesson that per-axis band-pass filtering \emph{before}
magnitude flattens the $\FI$ ratio was found precisely because the pipeline was
transparent enough to inspect.)
\item[Footprint and honesty.] The $\FI$ is two band-power sums and a ratio, well
inside the SAMD21's budget, and it runs untethered with no training set to drift away
from. The forest and CNN are research instruments that \emph{justify} and stress-test
the shipped rule; the rule itself is what a clinician can audit at the bedside.
\end{description}

\begin{keyidea}
Interpretability is not a consolation prize for a weaker model --- it is the headline
feature of a medical wearable. Cadence treats the CNN and RandomForest as
\emph{evidence} that its simple, transparent Freeze Index is sound, then ships the box
everyone can see straight through.
\end{keyidea}
